The normalized fourth/second moment ratio tends to $3$ in the maintained split
double average below $\theta<1/3$.  This is a fourth-moment theorem, not a CLT.

The pure off/off family contains a coherent reducible fibre.  For generic
$\lambda=(h_1,h_2)$ it occurs at
\[
t_\lambda=\frac{2}{c(h_1)-c(h_2)}
\]
and has exactly $2p-8$ valid affine points.  After extracting the spike,
Igusa's arithmetic open-image theorem \cite{BLV2009} gives geometric
irreducibility: the constants quotient is procyclic, so the geometric image
contains the commutator of an open arithmetic image, hence an open subgroup
of $\mathrm{SL}_2(\mathbf Z_\ell)$ for $\ell\ne p$.
The normalized elliptic local systems are pure of weight zero on their smooth
lisse loci; their multiplicative middle-extension stalks instead have
normalized eigenvalues $\pm p^{-1/2}$.
Apply the Fouvry--Kowalski--Michel trace-correlation argument \cite{FKM2013}
on the common smooth open and restore the boundedly many omitted points using
their bounded normalized traces.  This gives square-root quasi-orthogonality
between non-isomorphic geometric parameter classes.

\begin{theorem}[Corrected off/off variance]\label{thm:offoff}
\begin{equation}\label{eq:offoff}
\Var_{t\ne0}O(t)
\ll pE_+(A)+K^2p+K^4p^{1/2}+K^4,
\end{equation}
with an absolute implied constant.
\end{theorem}

Thus the off/off layer has collapsed onto the same pointwise Cayley-energy
bottleneck as the mixed layer.

\section{Exact trace-to-local-statistics bridge}

Let $U$ be an $N\times N$ unitary matrix with eigenvalues
$e^{2\pi i\theta_1},\ldots,e^{2\pi i\theta_N}$.  For a zero-sum integer tuple
$m_1+\cdots+m_k=0$, define
\[
D_U(m_1,\ldots,m_k)
 =\sum_{j_1,\ldots,j_k\ \mathrm{distinct}}
 e^{2\pi i\sum_{a=1}^km_a\theta_{j_a}}.
\]
For a set partition $\pi$ of $\{1,\ldots,k\}$, put
\[
\mu(\pi)=\prod_{B\in\pi}(-1)^{|B|-1}(|B|-1)!.
\]

\begin{theorem}[Factorial trace identity]\label{thm:factorial}
\begin{equation}\label{eq:factorial}
D_U(m_1,\ldots,m_k)
 =\sum_{\pi}\mu(\pi)
   \prod_{B\in\pi}
   \operatorname{Tr}U^{\sum_{a\in B}m_a}.
\end{equation}
\end{theorem}

This is M\"obius inversion on the partition lattice.  Let
$\phi\in\mathcal S(\mathbf R^{k-1})$ and periodize it at the mean-spacing scale.
Poisson summation gives the exact microscopic statistic
\begin{equation}\label{eq:trace-local}
C_{k,N}(U;\phi)
 =N^{-k}\sum_{m\in\Z^{k-1}}
  \widehat\phi(m/N)D_U(m_1,\ldots,m_k),
\end{equation}
where $m_k=-(m_1+\cdots+m_{k-1})$.  In particular,
\begin{equation}\label{eq:pair-local}
C_{2,N}(U;\phi)
 =N^{-2}\sum_{m\in\Z}\widehat\phi(m/N)
  \bigl(|\operatorname{Tr}U^m|^2-N\bigr).
\end{equation}
If $\widehat\phi$ has compact support, the relevant powers satisfy
$|m|=\Theta(N)$: microscopic local statistics probe the Heisenberg scale.

A weighted comparison of the quantities $D_U(m)$ with the matching random
matrix ensemble at every compact Fourier support is sufficient for the
corresponding band-limited microscopic correlations. Under positive normalized
averaging, the next theorem removes the band-limit. The required joint-trace
estimates for the correct desymmetrized cat-map block remain open.

\subsection{Positive correlation measures remove the band-limit}
Write $d=k-1$, unfold the circle to length $N$, and let $x_j=N\theta_j$.
The lifted anchored factorial measure is
\[
 \mu_N=\frac1N\sum_{j_1,\ldots,j_k\ \mathrm{distinct}}
 \sum_{\ell\in\mathbb Z^d}
 \delta_{(x_{j_1}-x_{j_k}+N\ell_1,\ldots,x_{j_d}-x_{j_k}+N\ell_d)}.
\]
It is positive, even when eigenvalues coincide, and
$\mu_N(\phi)=C_{k,N}(U;\phi)$.

\begin{lemma}[Uniform translated-cube bound]
\label{lem:factorial-cube-bound}
For $N>2$ and every $a\in\mathbb R^d$,
\[
 \mu_N(a+[0,1)^d)\le 2k^k\bigl(1+\mu_N([-2,2]^d)\bigr).
\]
The same bound holds under positive normalized averaging.
\end{lemma}
\begin{proof}
Let $n(t)$ count the periodically lifted points in $[t,t+2)$, with
multiplicity, and put $a_k=0$. For each tuple in the cube, integrate starting
points $t$ in the unit interval immediately preceding its anchor. Each of its
other points lies in $[t+a_i,t+a_i+2)$ apart from irrelevant endpoints.
Since $2<N$, every window contains at most one lift of each labelled point.
Integration over a period and then H\"older's inequality give
\[
 \mu_N(a+[0,1)^d)
 \le\frac1N\int_0^N\prod_{i=1}^k n(t+a_i)\,dt
 \le\frac1N\int_0^N n(t)^k\,dt.
\]
For the falling factorial $(r)_k$ one has
$r^k\le k^k((r)_k+r)$ for nonnegative integers $r$.
Indeed, for $r<k$ use $r^k\le k^{k-1}r$; for $r\ge k$ each factor in
$(r)_k$ is at least $r/k$. Moreover,
\[
 \frac1N\int_0^N n(t)\,dt=2,
\]
\[
 \frac1N\int_0^N(n(t))_k\,dt
 =\int\bigl(2-\operatorname{span}(0,z_1,\ldots,z_d)\bigr)_+\,d\mu_N(z)
 \le2\mu_N([-2,2]^d).
\]
The identity counts the possible starts of a common length-two window for
each distinct-label tuple; their total length is the displayed positive part.
This proves the bound. Positivity and normalization preserve it on averaging.
\end{proof}

\begin{theorem}[Extension from all compact Fourier supports]
\label{thm:factorial-schwartz-extension}
Fix $k\ge2$. Let $\mu_N,\nu_N$ be positive normalized averages of lifted
$k$-point factorial measures of $N$-point spectra. Suppose that $\nu_N$ has a
limit on all Schwartz tests and
\[
 \mu_N(f)-\nu_N(f)\longrightarrow0
\]
for every Schwartz $f$ with compact Fourier support, with no common upper
bound on that support. Then the same comparison holds for every Schwartz $f$.
\end{theorem}
\begin{proof}
Choose a nonzero nonnegative smooth compactly supported Fourier bump $\eta$,
and put $h=\mathcal F^{-1}\eta$. Since $h(0)>0$, for sufficiently small
$\epsilon>0$ the nonnegative band-limited Schwartz function
\[
 B(z)=4|h(\epsilon z)|^2/h(0)^2
\]
majorizes one on $[-2,2]^d$. The hypotheses bound $\nu_N(B)$ and
$\mu_N(B)$. Lemma~\ref{lem:factorial-cube-bound} therefore bounds both
measures on every translated unit cube by one constant $M$.
For $p>d$ and every Schwartz $g$ this implies
\[
 |\mu_N(g)|+|\nu_N(g)|
 \le C_{d,p}M\sup_z(1+|z|)^p|g(z)|.
\]
For a smooth compactly supported $\chi$ equal to one near zero, set
$f_R=\mathcal F^{-1}(\chi(\cdot/R)\widehat f)$.
Then $f_R\to f$ in Schwartz topology. The last inequality makes the two
approximation errors uniform in $N$. First send $N\to\infty$ at fixed $R$,
then $R\to\infty$.
\end{proof}

Thus a separate absolute high-frequency trace-tail estimate is unnecessary
for this positive-average, all-bandwidth criterion. This does not establish
such an absolute estimate, or extend a comparison confined to one fixed cap.
For example, for even $N$, the unit clock and a period-two configuration with
basis $\{0,a\}$, $0<a<2$, $a\ne1$, have pair measures
\[
 \sum_{r\in\mathbb Z}\delta_r-\sum_{r\in\mathbb Z}\delta_{Nr},
 \qquad
 \sum_{r\in\mathbb Z}\delta_{2r}
 +\frac12\sum_{r\in\mathbb Z}(\delta_{2r+a}+\delta_{2r-a})
 -\sum_{r\in\mathbb Z}\delta_{Nr}.
\]
They agree on tests supported in Fourier space inside $(-1/2,1/2)$ but
have different full limits. The theorem also requires positivity of the
averages. It addresses correlation tests, not a spacing-distribution theorem
without the additional point-process hypotheses needed for that conclusion.
